\section{缺席审判程序}
\begin{enumerate}
  \item 几乎每年都会考一道题。
  \item 核心。对席是常态，缺席是例外，所以\zd{尽可能保障被告人权利}。
  \item 检察院。
        \begin{enumerate}
          \item \zd{贪贿只要在境外}。依次严格。
          \item \zd{危恐境外及时审，还要最高检核准}。
          \item \zd{病半年}。严重疾病（丧失受审能力）中止超 6 月，被告人及其法代、近亲属申请 \hz 同意。
          \item \zd{已死亡}。
          \item 到案处理。
                \begin{enumerate}
                  \item 审查起诉 \yjt 重新审查。
                  \item 报请核准中 \yjt 撤回报请，重新审查。
                  \item 审理中。
                \end{enumerate}
        \end{enumerate}
  \item 法院。
        \begin{enumerate}
          \item 逃境外的。（犯罪地、被告人离境前的居住地 \hz 最高院指定）的中院。合议庭。
          \item 定罪标准不降低（事实清楚、证据充分）。审着发现罪名不符合案件范围，终止审理。
          \item 知情权。传票、起诉书副本域外送达 \yjt 收到不来才能缺席审判。
          \item 强制辩护。没有辩护人，则应当委托法援机构。
          \item 死亡处理。
                \begin{enumerate}
                  \item 一审、二审。终止审理；如果无罪，缺席还人清白。
                  \item 再审。缺席判决，依法做出裁判。
                \end{enumerate}
          \item 到案处理。
                \begin{enumerate}
                  \item 审理中 \yjt 重新审理。
                  \item 裁判生效后/交付执行前 \yjt 告知有权提出异议 \yjt 提 \yjt 应当重新审理。
                \end{enumerate}
        \end{enumerate}
  \item 当事人。
        \begin{enumerate}
          \item \zd{被告人或近亲属独立上诉权}。
          \item 辩护人经被告人同意上诉。
          \item 近亲属参加诉讼权。发表意见、出示证据、申请法庭通知证人/鉴定人出庭、进行辩论。
        \end{enumerate}
\end{enumerate}
